\section{两个模的张量积相对于正合列的性质}

记号约定:$A$是环,$E,E_{1},E_{2},E_{3}$是右$A$-模,$F,F_{1},F_{2},F_{3}$是左$A$-模.

\begin{proposition}{右模的张量积函子是右正合的}{}
	若序列$E_{1}\xlongrightarrow{f}E_{2}\xlongrightarrow{g}E_{3}\xlongrightarrow{}0$正合,则$E_{1}\otimes F\xlongrightarrow{f\otimes \id_{F}}E_{2}\otimes F\xlongrightarrow{g\otimes \id_{F}}E_{3}\otimes F\xlongrightarrow{}0$正合.
\end{proposition}

\begin{corollary}{左模的张量积函子是右正合的}{}
	若序列$F_{1}\xlongrightarrow{f}F_{2}\xlongrightarrow{g}F_{3}\xlongrightarrow{}0$正合,则$E\otimes F_{1}\xlongrightarrow{\id_{E}\otimes f}E\otimes F_{2}\xlongrightarrow{\id_{E}\otimes g}E\otimes F_{3}\xlongrightarrow{}0$正合.
\end{corollary}

\begin{remark}{张量积函子不一定是左整合的}{}
	若序列$0\xlongrightarrow{}E_{1}\xlongrightarrow{f}E_{2}\xlongrightarrow{g}E_{3}\xlongrightarrow{}0$正合,则$0\xlongrightarrow{}E\otimes F_{1}\xlongrightarrow{\id_{E}\otimes f}E\otimes F_{2}\xlongrightarrow{\id_{E}\otimes g}E\otimes F_{3}\xlongrightarrow{}0$未必正合.例如,若
	\begin{align*}
		A=\Z,E_{1}=2\Z,E_{2}=\Z,F_{2}=\Z/2\Z,
	\end{align*}
	则$\im\left(f\otimes\id_{F}\right)=0$.因此,计算$x\otimes y$时,应该明确它属于$E_{1}\otimes_{A}F$还是属于$E_{2}\otimes_{A}F$.
\end{remark}

\begin{proposition}{满射的线性映射的张量积的核}{}
	若序列$E_{1}\xlongrightarrow{f_{1}}E_{2}\xlongrightarrow{f_{2}}E_{3}\xlongrightarrow{}0$和$F_{1}\xlongrightarrow{g_{1}}F_{2}\xlongrightarrow{g_{2}}F_{3}\xlongrightarrow{}0$都正合,则$f_{2}\otimes g_{2}$是满射,且
	\begin{align*}
		\ker\left(f_{2}\otimes g_{2}\right)=\im\left(f_{1}\otimes\id_{F_{2}}\right)+\im\left(\id_{E_{2}}\otimes g_{1}\right).
	\end{align*}
\end{proposition}

\begin{corollary}{商模的张量积同构}{}
	设$E_{1}$是$E_{2}$的子模,$F_{1}$是$F_{2}$的子模,$\iota_{1}:E_{1}\otimes F_{2}\hookrightarrow E_{2}\otimes F_{2}$和$\iota_{2}:E_{2}\otimes F_{1}\hookrightarrow E_{2}\otimes F_{2}$是典范嵌入,则有典范的$\Z$-模同构
	\begin{align*}
		\left(E_{2}/E_{1}\right)\otimes_{A}\left(F_{2}/F_{1}\right)\simeq\left(E_{2}\otimes_{A}F_{2}\right)/\left(\im\left(\iota_{1}\right)+\im\left(\iota_{2}\right)\right).
	\end{align*}
\end{corollary}

\begin{remark}{商多模的张量积同构}{}
	上述推论可以推广到多模,只需要做如下替换:
	\begin{itemize}
		\item $E_{1}$是右$A$-模$\mapsto$ $E$是$\left(\mathsf{B};A,\mathsf{C}\right)$-多模;
		\item $F$是左$A$-模$\mapsto$ $F$是$\left(A,\mathsf{B}^{\prime};\mathsf{C}^{\prime}\right)$-多模;
		\item $E_{1}$是$E_{2}$的右$A$-子模$\mapsto$ $E_{1}$是$E_{2}$的$\left(\mathsf{B};A,\mathsf{C}\right)$-子多模;
		\item $F_{1}$是$F_{2}$的左$A$-子模$\mapsto$ $F_{1}$是$F_{2}$的$\left(A,\mathsf{B}^{\prime};\mathsf{C}^{\prime}\right)$-子多模;
		\item 典范的$\Z$-模同构$\mapsto$ 典范的$\left(\mathsf{B},\mathsf{B}^{\prime};\mathsf{C},\mathsf{C}^{\prime}\right)$.
	\end{itemize}
	其中,$\mathsf{B},\mathsf{B}^{\prime},\mathsf{C},\mathsf{C}^{\prime}$是四族环.
\end{remark}

\begin{corollary}{}{}
	设$\mathfrak{a}$是$A$的左理想,则有典范$\Z$-模同构$\left(A/\mathfrak{a}\right)\otimes_{A}F\simeq F/\mathfrak{a}F$.
\end{corollary}

\begin{remark}{}{}
	特别地,当$A=\Z,n\in\Z$时,$\left(\Z/n\Z\right)\otimes_{\Z}E$典范同构于$F/nF$.
\end{remark}

\begin{corollary}{}{}
	设$A$是交换环,$\mathfrak{a}$是$A$的理想,$\mathfrak{a}F=0$,则有典范$A/\mathfrak{a}$-模同构$E\otimes F\simeq\left(E/\mathfrak{a}E\right)\otimes_{A/\mathfrak{a}}F$.
\end{corollary}

\begin{corollary}{}{}
	设$A$是交换环,$\mathfrak{a},\mathfrak{b}$是$A$的理想,则有典范的$A$-模同构$\left(_{d}A_{s}\right)/\mathfrak{a}\otimes_{A}\left(_{d}A_{s}\right)/\mathfrak{b}\simeq\left(_{d}A_{s}\right)/\left(\mathfrak{a}+\mathfrak{b}\right)$.
\end{corollary}
